% Normal-approximations-to-binomial.tex
% Topic: Normal Approximations to a Binomial

When $n$ is large, the binomial distribution $B(n,p)$ can be approximated by a normal distribution. Specifically,
\[
X \sim B(n,p) \approx N\!\left(\mu, \sigma^2\right), \qquad \mu = np, \quad \sigma^2 = np(1-p).
\]

\textbf{When is the approximation valid?} A common guideline is that both $np \geq 5$ and $n(1-p) \geq 5$.

\textbf{Continuity correction.} Because the binomial is discrete and the normal is continuous, a continuity correction improves accuracy:
\[
P(X = k) \approx P\!\left(k - \tfrac{1}{2} \leq X \leq k + \tfrac{1}{2}\right).
\]
For example:
\[
P(X \leq k) \approx P\!\left(X \leq k + \tfrac{1}{2}\right) \quad \text{(normal approximation)}.
\]

\subsubsection*{Illustration ($p=\tfrac{1}{2}$, $\mu=np$, $\sigma^2=np(1-p)$)}
We fix $p=\frac{1}{2}$ so the binomial is symmetric---the histogram of $P(X=k)$ is easier to compare to the bell-shaped curve. For each $n$, the matching normal density is $N(np,\ np(1-p))$ (plotted as a curve on the same $k$-axis). Only for larger $n$ does the guideline $np\ge 5$, $n(1-p)\ge 5$ hold; the diagrams show visually why the approximation is guarded for small~$n$.

\begin{center}\small
\begin{tikzpicture}
\begin{groupplot}[
    group style={
        group size=1 by 3,
        vertical sep=48pt,
        y descriptions at=edge left,
    },
    width=0.88\textwidth,
    height=3.55cm,
    axis lines=left,
    enlarge x limits=0.06,
    ymin=0,
    grid=both,
    major grid style={gray!22},
    tick label style={font=\footnotesize},
    label style={font=\small},
    title style={font=\normalsize\bfseries, yshift=2pt},
    every axis plot/.append style={semithick},
    clip mode=individual,
    legend style={
        font=\footnotesize,
        anchor=south,
        legend columns=-1,
        column sep=0.55em,
    },
    legend to name=binomnormalapproachlegend,
]

% --- Top: n = 2 ---
\nextgroupplot[
    title={$n=2$;\quad $\mu=1$, $\sigma^2=\tfrac12$},
    xmin=-0.45, xmax=2.45,
    ymax=0.58,
    xtick={0,1,2},
    xticklabels={,,},
    ytick={0,0.1,0.2,0.3,0.4,0.5},
    scaled ticks=false,
    ylabel={$P(X=k)$},
]

\addplot+[ycomb, thick, mark=*, mark size=1.35pt, blue!75!black]
coordinates {(0,0.25)(1,0.5)(2,0.25)};
\addplot+[thick, samples=140, smooth, blue!65!cyan, dashed, domain=-0.45:2.45]
{1/(sqrt(0.5)*sqrt(2*pi))*exp(-(x-1)^2/(2*0.5))};
\addlegendentry{binomial pmf}
\addlegendentry{$N(np,npq)$ pdf}

% --- Middle: n = 5 ---
\nextgroupplot[
    title={$n=5$;\quad $\mu=\tfrac52$, $\sigma^2=\tfrac54$},
    xmin=-0.45, xmax=5.45,
    ymax=0.38,
    xtick={0,1,...,5},
    xticklabels={,,,,, },
    ytick={0,0.1,0.2,0.3},
    ylabel={$P(X=k)$},
]

\addplot+[ycomb, thick, mark=*, mark size=1.2pt, green!55!black, forget plot]
coordinates {
(0,0.03125)(1,0.15625)(2,0.3125)(3,0.3125)(4,0.15625)(5,0.03125)
};
\addplot+[thick, samples=160, smooth, green!45!gray, dashed, domain=-0.45:5.45, forget plot]
{1/(sqrt(1.25)*sqrt(2*pi))*exp(-(x-2.5)^2/(2*1.25))};

% --- Bottom: n = 20 ---
\nextgroupplot[
    title={$n=20$;\quad $\mu=10$, $\sigma^2=5$},
    xmin=-0.6, xmax=20.6,
    ymax=0.20,
    xtick={0,5,...,20},
    ytick={0,0.05,0.10,0.15,0.20},
    xlabel={successes $k$},
    ylabel={$P(X=k)$},
]

\addplot+[ycomb, thick, mark=*, mark size=0.9pt, red!72!black, forget plot]
coordinates {
(0,0.000001)(1,0.000019)(2,0.000181)(3,0.001087)(4,0.004621)
(5,0.014786)(6,0.036964)(7,0.073929)(8,0.120134)(9,0.160179)
(10,0.176197)(11,0.160179)(12,0.120134)(13,0.073929)(14,0.036964)
(15,0.014786)(16,0.004621)(17,0.001087)(18,0.000181)(19,0.000019)
(20,0.000001)
};
\addplot+[thick, samples=220, smooth, red!40!orange, dashed, domain=-0.6:20.6, forget plot]
{1/(sqrt(5)*sqrt(2*pi))*exp(-(x-10)^2/(2*5))};

\end{groupplot}
\end{tikzpicture}
\vspace{0.65em}
\pgfplotslegendfromname{binomnormalapproachlegend}
\vspace{0.45em}
\noindent\footnotesize\textit{Increasing $n$ (reading downward): the pmf becomes more bell-shaped and tracks the $N(np,npq)$ curve more closely. Here $p=\frac12$ so $pq=\frac14$ and $\mathrm{Var}(X)=npq$.}
\end{center}

\subsubsection*{Continuity correction in one picture}
Discrete pmf masses at integers sit under a smooth $N(np,npq)$ curve; approximating $P(X=k)$ uses a width-one slice $[k-\tfrac12,k+\tfrac12]$.

\begin{center}\small
\begin{tikzpicture}
  \begin{axis}[
      DLfigaxis,
      width=0.94\textwidth,
      height=4.05cm,
      xmin=4.55,
      xmax=13.95,
      ymin=-0.001,
      ymax=0.28,
      xlabel={count of successes $k$ (illustrative $B(16,\tfrac12)$)},
    ]

    \def\binpdf{exp(-((x-8)^2)/(2*4))/sqrt(2*pi*4)}

    \addplot[
      orange!93!black, draw=none, fill=orange!93!black, fill opacity=0.25,
      samples=340, smooth, domain=7.5:8.5,
    ]
      {\binpdf}\closedcycle;

    \addplot+[ycomb, thick, mark size=2pt,
      purple!74!black, forget plot]
    coordinates {
      (8,0.177)(7,0.158)(9,0.158)(6,0.104)(10,0.104)(5,0.047)(11,0.047)(4,0.016)(12,0.016)
    };

    \addplot[thick, blue!73!cyan, densely dashed, samples=360, smooth, domain=4.6:12.4]
      {\binpdf};

    \draw[|-|, thick] (axis cs:7.5,0.21) -- (axis cs:8.5,0.21)
      node[midway, above=2pt, font=\scriptsize] {$k-\tfrac12$ to $k+\tfrac12$};

  \end{axis}
\end{tikzpicture}\\[0.25em]

\noindent\footnotesize\textit{Continuity correction: approximate the mass at integer $k$ using the normal density on a width-one interval around $k$. Purple spikes: exact binomial pmf; dashed curve: $N(np,np(1-p))$.}
\end{center}

